\chapter{引用、指针和非拥有视图}

\section{问题入口：函数到底改没改原对象}

很多 \cpp 初学者第一次写函数时，会遇到一个反直觉问题：

\begin{lstlisting}[language=C++,caption={按值传参不会修改原对象}]
void add_bonus(int score)
{
    score += 5;
}

int main()
{
    int value = 90;
    add_bonus(value);
    std::cout << value << '\n'; // 仍然是 90
}
\end{lstlisting}

这段代码没有错。函数参数 \code{score} 是一个新对象，它从 \code{value} 复制得到。修改 \code{score} 只会修改副本。这个反例推出第一条规则：看到函数参数时，先判断它是在复制对象、引用对象，还是通过指针间接访问对象。

如果函数确实要修改调用者传入的对象，可以用引用：

\begin{lstlisting}[language=C++,caption={非 const 引用表达可修改}]
void add_bonus(int& score)
{
    score += 5;
}
\end{lstlisting}

\code{int&} 不是一个新的整数对象，而是已有整数对象的别名。调用 \code{add_bonus(value)} 后，函数内部的 \code{score} 和外部的 \code{value} 指向同一个对象。

\section{引用的核心：必须绑定到有效对象}

引用一旦创建，就必须绑定到某个对象。它不能“空着”，也不能随便改绑到另一个对象：

\begin{lstlisting}[language=C++,caption={引用绑定后就是别名}]
int first = 1;
int second = 2;
int& ref = first;
ref = second; // 不是改绑，而是把 first 改成 2
\end{lstlisting}

执行后 \code{first} 变成 \code{2}，\code{ref} 仍然是 \code{first} 的别名。这个细节很重要，因为它解释了为什么引用适合表达“这个参数一定存在”。如果对象可能不存在，引用就不是合适表达。

\section{指针表达可空和可重新指向}

指针保存对象地址，可以为空，也可以重新指向别的对象：

\begin{lstlisting}[language=C++,caption={指针可以表示可选对象}]
void print_score(const int* score)
{
    if (score == nullptr) {
        std::cout << "missing\n";
        return;
    }
    std::cout << *score << '\n';
}
\end{lstlisting}

\code{*score} 表示访问指针指向的对象。访问前必须确认它不是 \code{nullptr}。这推出第二条规则：接口如果用裸指针表达观察关系，就必须清楚说明它是否允许为空。允许为空时，函数内部要检查；不允许为空时，优先考虑引用。

\section{不要返回局部对象的地址}

指针和引用都不拥有对象，所以最常见错误是指向已经销毁的局部变量：

\begin{lstlisting}[language=C++,caption={返回局部变量地址是错误}]
const int* bad_pointer()
{
    int value = 42;
    return &value; // 函数结束后 value 已经销毁
}
\end{lstlisting}

函数返回后，\code{value} 的生命周期结束。调用者拿到的地址已经悬空。正确做法通常是返回值：

\begin{lstlisting}[language=C++,caption={返回值让调用者拥有结果}]
int make_value()
{
    return 42;
}
\end{lstlisting}

现代 \cpp 返回值并不意味着低效。编译器会进行返回值优化，必要时也可以移动对象。为了省一次可能不存在的复制而返回悬空引用，是非常不划算的。

\section{span：一段连续数据的非拥有视图}

当函数需要观察数组、\code{vector} 或 \code{array} 的一段连续元素时，\code{std::span} 比裸指针加长度更清楚：

\begin{lstlisting}[language=C++,caption={用 span 表达连续区间}]
int sum_scores(std::span<const int> scores)
{
    int total = 0;
    for (int score : scores) {
        total += score;
    }
    return total;
}
\end{lstlisting}

\code{std::span<const int>} 表达三件事：它不拥有元素；它只读；它知道长度。调用者可以传 \code{std::vector<int>}、\code{std::array<int, N>} 或普通数组。函数内部不用同时维护 \code{int* data} 和 \code{std::size_t size} 两个参数，也减少了长度不匹配的风险。

\section{string\_view：字符串内容的非拥有视图}

\code{std::string_view} 类似 \code{span<const char>}，但专门表示字符串内容。它适合只观察文本：

\begin{lstlisting}[language=C++,caption={字符串视图适合只读参数}]
bool starts_with(std::string_view text, std::string_view prefix)
{
    return text.size() >= prefix.size() &&
           text.substr(0, prefix.size()) == prefix;
}
\end{lstlisting}

它不分配内存，也不要求调用者一定传 \code{std::string}。字符串字面量、\code{std::string} 和子串都能传入。但它不拥有内容，所以不能长期保存一个指向临时字符串的 \code{string_view}：

\begin{lstlisting}[language=C++,caption={保存临时字符串视图会悬空}]
std::string_view bad_view()
{
    std::string text = "temporary";
    return std::string_view{text}; // text 离开函数后销毁
}
\end{lstlisting}

非拥有视图的好处是轻，代价是生命周期要求更严格。只要你把视图保存到函数调用之后，就必须能证明底层对象活得更久。

\section{接口选择表}

\begin{longtable}{p{0.28\linewidth}p{0.28\linewidth}p{0.34\linewidth}}
\toprule
意图 & 推荐形态 & 说明 \\
\midrule
复制一个便宜值 & \code{int value} & 标量和小对象可以按值 \\
只观察一个大对象 & \code{const T&} & 不复制，不修改 \\
修改一个必定存在的对象 & \code{T&} & 调用点能看出会修改 \\
观察可能不存在的对象 & \code{const T*} & \code{nullptr} 表示没有 \\
接管独占所有权 & \code{std::unique_ptr<T>} & 调用点需要 \code{std::move} \\
观察连续数据 & \code{std::span<const T>} & 不拥有，带长度 \\
观察文本 & \code{std::string_view} & 不拥有字符内容 \\
\bottomrule
\end{longtable}

\section{完整案例：校验用户输入}

假设要校验用户名：不能为空，只能包含字母、数字和下划线，长度不能超过 \code{32}。这个函数只观察文本，不保存文本，所以用 \code{std::string_view}：

\begin{lstlisting}[language=C++,caption={用户名校验}]
#include <cctype>
#include <string_view>

constexpr std::size_t USERNAME_MAX_LENGTH = 32;

bool is_valid_username(std::string_view name)
{
    if (name.empty() || name.size() > USERNAME_MAX_LENGTH) {
        return false;
    }
    for (unsigned char ch : name) {
        if (std::isalnum(ch) || ch == '_') {
            continue;
        }
        return false;
    }
    return true;
}
\end{lstlisting}

这里 \code{USERNAME_MAX_LENGTH} 用常量命名，而不是把 \code{32} 散在代码里。循环里用 \code{unsigned char} 是因为 \code{std::isalnum} 对普通负 \code{char} 的行为不安全。一个小函数也能体现类型、生命周期、常量和边界处理。

\begin{exercisebox}
实现 \code{trim(std::string_view text)}，返回去掉首尾空白后的 \code{std::string_view}。要求不分配新字符串。写测试覆盖空字符串、全空白、无空白、只有左空白、只有右空白。然后回答：返回的视图依赖谁的生命周期？
\end{exercisebox}
